\clearpage
\section{Porównanie wyników i analiza}

\subsection{Synteza wyników badań}

\subsubsection{Zestawienie wyników testów wydajności}

Na~podstawie przeprowadzonych testów wydajności (rozdział~\ref{sec:testy-wydajnosci}) można stwierdzić, że~oba silniki wykazują odmienne charakterystyki wydajnościowe w~zależności od~scenariusza testowego.

\begin{table}[h!]
    \centering
    \caption{Zestawienie zbiorcze wyników wydajnościowych}
    \label{tab:performance-summary}
    \begin{tabular}{|l|c|c|}
        \hline
        \textbf{Kryterium} & \textbf{Unity} & \textbf{Unreal Engine} \\
        \hline\hline
        Średni FPS (scenariusz niski) & [DATA] & [DATA] \\
        \hline
        Średni FPS (scenariusz średni) & [DATA] & [DATA] \\
        \hline
        Średni FPS (scenariusz wysoki) & [DATA] & [DATA] \\
        \hline
        Średnie zużycie VRAM & [DATA] MB & [DATA] MB \\
        \hline
        Średnia liczba draw calls & [DATA] & [DATA] \\
        \hline
        Stabilność frame time & [DATA] & [DATA] \\
        \hline
    \end{tabular}
\end{table}

\subsubsection{Zestawienie analizy funkcjonalności}

Analiza możliwości (rozdział~6) wykazała, że~oba silniki oferują bogate zestawy funkcjonalności, jednak~skierowane do~nieco odmiennych grup docelowych.

\subsection{Analiza wielokryterialna}

\subsubsection{Macierz porównawcza}

\begin{table}[h!]
    \centering
    \caption{Macierz wielokryterialnego porównania silników gier}
    \label{tab:comparison-matrix}
    \begin{tabular}{|l|c|c|p{3cm}|}
        \hline
        \textbf{Kryterium} & \textbf{Unity} & \textbf{Unreal} & \textbf{Uwagi} \\
        \hline\hline
        Wydajność 2D & ★★★★★ & ★★★☆☆ & Unity zoptymalizowane pod~2D \\
        \hline
        Wydajność 3D & ★★★★☆ & ★★★★★ & Unreal lepsze w~AAA 3D \\
        \hline
        Jakość grafiki & ★★★★☆ & ★★★★★ & Unreal oferuje Nanite, Lumen \\
        \hline
        Łatwość nauki & ★★★★★ & ★★★☆☆ & Unity bardziej przystępne \\
        \hline
        Dokumentacja & ★★★★★ & ★★★☆☆ & Unity ma~lepszą dokumentację \\
        \hline
        Wsparcie mobilne & ★★★★★ & ★★★☆☆ & Unity dominuje na~mobile \\
        \hline
        Społeczność & ★★★★★ & ★★★★☆ & Większa społeczność Unity \\
        \hline
        Asset Store & ★★★★★ & ★★★★☆ & Więcej zasobów dla~Unity \\
        \hline
        Blueprint/Visual & ★★★☆☆ & ★★★★★ & Blueprints bardziej zaawansowane \\
        \hline
        Kod C\# vs C++ & ★★★★☆ & ★★★★☆ & C\# łatwiejszy, C++ wydajniejszy \\
        \hline
        Licensing & ★★★★★ & ★★★★★ & Oba bezpłatne dla~indie \\
        \hline
    \end{tabular}
\end{table}

\subsubsection{Analiza wag kryteriów}

Znaczenie poszczególnych kryteriów różni się w~zależności od~typu projektu:

\begin{itemize}
    \item \textbf{Gry indie} -- priorytet: łatwość nauki, koszt, społeczność
    \item \textbf{Gry mobilne} -- priorytet: wydajność, optymalizacja, wsparcie platform
    \item \textbf{Gry AAA} -- priorytet: jakość grafiki, zaawansowane funkcje, skalowalność
    \item \textbf{Gry edukacyjne} -- priorytet: prostota, dokumentacja, stabilność
\end{itemize}

\subsection{Przypadki użycia}

\subsubsection{Gry indie}

\paragraph{Rekomendacja}: Unity

\paragraph{Uzasadnienie}:
\begin{itemize}
    \item Niższy próg wejścia dla~początkujących deweloperów
    \item Bogaty Asset Store z~dostępnymi cenowo zasobami
    \item Większa społeczność -- łatwiej znaleźć pomoc
    \item Szybsze prototypowanie
    \item Mniejsze wymagania sprzętowe dla~deweloperów
\end{itemize}

\paragraph{Wyjątki}:
Jeśli gra wymaga grafiki najwyższej jakości (photorealistic), rozważ Unreal Engine.

\subsubsection{Gry mobilne}

\paragraph{Rekomendacja}: Unity

\paragraph{Uzasadnienie}:
\begin{itemize}
    \item Lepsza optymalizacja pod~platformy mobilne
    \item Mniejsze rozmiary buildu
    \item Lepsze wsparcie dla~starszych urządzeń
    \item Więcej narzędzi i~assetów mobilnych
    \item Większość gier mobilnych używa Unity (udowodniona skuteczność)
\end{itemize}

\paragraph{Statystyki}:
Według danych z~2023 roku, około 70\% gier mobilnych na~iOS i~Android zostało stworzonych w~Unity~\cite{statista_unity_market, unity_gaming_report}.

\subsubsection{Gry AAA}

\paragraph{Rekomendacja}: Unreal Engine

\paragraph{Uzasadnienie}:
\begin{itemize}
    \item Wyższa jakość grafiki out-of-the-box
    \item Nanite -- rendering miliardów poligonów
    \item Lumen -- dynamiczne global illumination
    \item Lepsze wsparcie dla~dużych zespołów
    \item Sprawdzone w~produkcjach AAA (Fortnite, Gears of~War)
\end{itemize}

\paragraph{Przykłady}:
Unreal Engine wykorzystywano w~produkcjach takich jak: Final Fantasy VII Remake, Jedi: Fallen Order, Borderlands 3.

\subsubsection{Gry VR/AR}

\paragraph{Rekomendacja}: Zależnie od~wymagań

\paragraph{Unity dla}:
\begin{itemize}
    \item Aplikacje edukacyjne VR/AR
    \item Mobilny AR (ARCore, ARKit)
    \item Projekty wymagające szybkiego rozwoju
\end{itemize}

\paragraph{Unreal dla}:
\begin{itemize}
    \item High-end VR experiences
    \item Architekturalna wizualizacja VR
    \item Training simulations wymagające fotorealizmu
\end{itemize}

\subsection{Weryfikacja hipotez badawczych}

Na~początku pracy (rozdział~\ref{sec:wstep}) postawiono następujące hipotezy badawcze:

\subsubsection{Hipoteza 1: Wydajność renderowania 3D}

\textbf{Hipoteza}: Unreal Engine oferuje lepszą wydajność renderowania złożonych scen 3D niż~Unity.

\textbf{Weryfikacja}: [POTWIERDZONA/ODRZUCONA -- wypełnij po~analizie danych]

\textbf{Uzasadnienie}: Na~podstawie testów wydajności (Tabela~\ref{tab:results-high}) zaobserwowano, że...

\subsubsection{Hipoteza 2: Łatwość nauki}

\textbf{Hipoteza}: Unity charakteryzuje się niższym progiem wejścia dla~początkujących deweloperów niż~Unreal Engine.

\textbf{Weryfikacja}: POTWIERDZONA

\textbf{Uzasadnienie}: Analiza wywiadów (rozdział~wywiady-analiza) wykazała, że~100\% respondentów z~doświadczeniem poniżej 2~lat oceniło Unity jako bardziej przystępne. Składają się na~to:
\begin{itemize}
    \item Lepiej udokumentowane API
    \item Większa dostępność tutoriali dla~początkujących
    \item C\# jako język bardziej przyjazny niż~C++
    \item Prostszy interfejs edytora
\end{itemize}

\subsubsection{Hipoteza 3: Wsparcie mobilne}

\textbf{Hipoteza}: Unity oferuje lepsze wsparcie i~optymalizację dla~platform mobilnych niż~Unreal Engine.

\textbf{Weryfikacja}: POTWIERDZONA

\textbf{Uzasadnienie}:
\begin{itemize}
    \item Mniejsze rozmiary buildów mobilnych w~Unity
    \item Lepsza optymalizacja dla~urządzeń niskiej klasy
    \item Większy ekosystem mobile-specific assetów
    \item Dominacja na~rynku gier mobilnych (70\% udziału)
\end{itemize}

\subsubsection{Hipoteza 4: Jakość grafiki}

\textbf{Hipoteza}: Unreal Engine umożliwia osiągnięcie wyższej jakości grafiki niż~Unity przy~porównywalnym nakładzie pracy.

\textbf{Weryfikacja}: POTWIERDZONA

\textbf{Uzasadnienie}:
\begin{itemize}
    \item Technologie Nanite i~Lumen oferują funkcje niedostępne w~Unity
    \item Lepsze domyślne materiały i~shadery
    \item Zaawansowane efekty post-processingu out-of-the-box
    \item Większość projektów wymagających fotorealizmu wykorzystuje Unreal
\end{itemize}

\subsection{Ograniczenia badań}

\subsubsection{Ograniczenia metodologiczne}

\begin{itemize}
    \item \textbf{Ograniczona liczba scenariuszy testowych} -- skupiono się na~grze typu bullet-hell, co~nie~pokrywa wszystkich możliwych zastosowań silników
    \item \textbf{Pojedyncza konfiguracja sprzętowa} -- testy przeprowadzono tylko na~jednym zestawie komputerowym
    \item \textbf{Mała próba wywiadów} -- 8~respondentów może nie~reprezentować całej społeczności deweloperów
\end{itemize}

\subsubsection{Ograniczenia techniczne}

\begin{itemize}
    \item \textbf{Wersje silników} -- wyniki dotyczą konkretnych wersji Unity i~Unreal; nowsze wersje mogą mieć odmienną wydajność
    \item \textbf{Wpływ object poolingu} -- optymalizacja wpływa na~wyniki; bez~niej różnice mogłyby być większe
    \item \textbf{Profilowanie} -- NVIDIA Nsight może wprowadzać własny narzut wydajnościowy
\end{itemize}

\subsubsection{Ograniczenia czasowe}

\begin{itemize}
    \item Silniki gier rozwijają się dynamicznie -- wyniki mogą dezaktualizować się w~ciągu roku
    \item Nie~testowano funkcji wprowadzonych w~najnowszych wersjach beta
\end{itemize}

\subsection{Implikacje praktyczne}

\subsubsection{Dla deweloperów indywidualnych}

\begin{itemize}
    \item Rozpoczynając naukę tworzenia gier, Unity stanowi bezpieczniejszy wybór
    \item Dla projektów 2D, Unity jest jednoznacznie lepszym wyborem
    \item Inwestycja w~naukę C++ może być wartościowa długoterminowo
\end{itemize}

\subsubsection{Dla małych zespołów (2-10 osób)}

\begin{itemize}
    \item Unity pozwala na~szybsze MVP i~iteracje
    \item Unreal wymaga co~najmniej jednego doświadczonego programisty C++
    \item Asset Store Unity oferuje więcej ready-to-use rozwiązań
\end{itemize}

\subsubsection{Dla studiów AAA}

\begin{itemize}
    \item Unreal Engine jest standardem przemysłowym dla~gier 3D wysokiej jakości
    \item Wsparcie Epic Games dla~dużych projektów jest lepsze
    \item Source code access w~Unreal daje większą kontrolę
\end{itemize}
